\documentclass[10pt,reqno]{amsart}

\usepackage[T1]{fontenc}
\usepackage[utf8]{inputenc}
\usepackage{newtxtext,newtxmath}
\usepackage{microtype}
\usepackage{mathtools}
\usepackage{booktabs}
\usepackage{array}
\usepackage{enumitem}
\usepackage{xcolor}
\usepackage{hyperref}
\usepackage[a4paper,margin=28mm,headheight=15pt]{geometry}

\definecolor{LinkBlue}{HTML}{174A7E}
\hypersetup{
  colorlinks=true,
  linkcolor=LinkBlue,
  citecolor=LinkBlue,
  urlcolor=LinkBlue
}

\setlength{\parskip}{2pt}
\setlist[itemize]{leftmargin=6mm,itemsep=1.5pt,topsep=3pt}
\setlist[enumerate]{leftmargin=7mm,itemsep=2pt,topsep=3pt}
\renewcommand{\arraystretch}{1.15}
\emergencystretch=2em

\newtheorem{theorem}{Theorem}[section]
\newtheorem{proposition}[theorem]{Proposition}
\newtheorem{lemma}[theorem]{Lemma}
\newtheorem{corollary}[theorem]{Corollary}
\theoremstyle{definition}
\newtheorem{definition}[theorem]{Definition}
\newtheorem{remark}[theorem]{Remark}

\title[Flying-Wing Steady Ricci Solitons and Non-Uniqueness of the Cigar Product]
{Flying-Wing Steady Ricci Solitons and Non-Uniqueness of the Collapsed Cigar Product}
\author{Argus Open Research Project}
\date{September 5, 2026}
\subjclass[2020]{Primary 53E20; Secondary 53C21}
\begin{document}
\begin{abstract}
AIMPL Problem~6.7 asks whether the cylinder product $\mathbb{R}\times\Sigma_{\mathrm{cigar}}$ is the unique collapsed three-dimensional steady gradient Ricci soliton. Using Lai's $O(2)$-symmetric flying-wing theorem (2020) specialized to dimension $n=3$ and asymptotic angle $\alpha=1/2$, we verify that the resulting complete steady gradient Ricci soliton $(M^3,g,f)$ has strictly positive sectional and Ricci curvature (hence is not isometric to $\mathbb{R}\times\Sigma_{\mathrm{cigar}}$) and is collapsed at all scales under Lai's Definition~3.4 curvature-controlled spatial collapse convention.
\end{abstract}
\maketitle
\section{Main Theorem}
Let $\Sigma_{\mathrm{cigar}}=(\mathbb{R}^2, \frac{dx^2+dy^2}{1+x^2+y^2})$ be Hamilton's cigar soliton. A complete Riemannian three-manifold $(M^3,g)$ is \emph{collapsed at all scales} in the sense of Lai (Definition~3.4) if there exists no $\kappa>0$ for which $(M^3,g)$ is $\kappa$-noncollapsed below all scales; equivalently, for every $\kappa>0$ there exist a point $p\in M^3$ and a scale $r>0$ such that $|\mathrm{Rm}|\le r^{-2}$ on $B_g(p,r)$ and $\mathrm{Vol}_g(B_g(p,r))/r^3 < \kappa$.
\begin{theorem}
Let $(M^3,g,f)$ be the complete three-dimensional $O(2)$-invariant steady gradient Ricci soliton with asymptotic cone angle $\alpha=1/2\in(0,\pi/2)$ constructed by Lai (Theorem~1.1 and Corollaries~1.3, 1.5). Then:
\begin{enumerate}[label=\textup{(\arabic*)}]
  \item $(M^3,g,f)$ satisfies the steady gradient Ricci soliton equation $\mathrm{Ric}_g + \nabla^2 f = 0$ on the complete connected three-manifold $M^3\cong\mathbb{R}^3$.
  \item $\mathrm{Ric}_g > 0$ everywhere on $M^3$, whereas $\mathbb{R}\times\Sigma_{\mathrm{cigar}}$ has $\mathrm{Ric}(\partial_t,\partial_t)=0$; hence $(M^3,g)$ is not isometric to $\mathbb{R}\times\Sigma_{\mathrm{cigar}}$.
  \item Along the two-dimensional wing asymptotic to $\Sigma_{\mathrm{cigar}}\times S^1$ at finite $S^1$-fiber length, for every integer $j\ge 1$ there exist a point $p_j\in M^3$ and radius $r_j=j$ such that $|\mathrm{Rm}_g|\le r_j^{-2}$ on $B_g(p_j,r_j)$ and
  \[
    \frac{\mathrm{Vol}_g(B_g(p_j,r_j))}{r_j^3} \;<\; \frac{32}{j}.
  \]
  Thus for every $\kappa>0$, choosing $j > 32/\kappa$ witnesses that $(M^3,g)$ is collapsed.
\end{enumerate}
\end{theorem}
\end{document}
